\chapter{Objetivos}
\label{cap:objetivos}

\section{Objetivo general}

Distinguir las consecuencias mecánicas de un AMM de producto constante sin comisiones de los efectos de reglas de agentes especificadas, y caracterizar el seguimiento con arbitraje acotado, la sensibilidad al retiro de liquidez y el desempeño de modelos reducidos dentro de un micromundo CEX/DEX sintético y reproducible.

\section{Objetivos específicos}

\begin{enumerate}[label=\textbf{OE\arabic*.}, leftmargin=3em]
  \item Especificar y verificar un simulador discreto de un AMM de producto constante con una referencia CEX exógena, reglas de ruido, tendencia, arbitraje y provisión de liquidez, y registrar sus ejecuciones de forma reproducible.
  \item Derivar la tasa de corrección ideal que implica la capacidad de arbitraje y compararla con trayectorias del simulador de horizonte finito bajo los flujos de órdenes configurados.
  \item Caracterizar cómo el disparador de retiro único existente cambia el seguimiento simulado con órdenes fijas y proporcionales a las reservas, incluida la ablación del agente de pánico, sin tratar la regla como comportamiento de proveedores estimado empíricamente.
  \item Evaluar la precisión descriptiva de mapas reducidos en trayectorias reservadas del simulador frente a la persistencia y un AR(1) escalar, comparar la reproducción recursiva sensible a la banda con un mapa afín y delimitar las afirmaciones predictivas de la clasificación de regímenes.
\end{enumerate}

\section{Correspondencia entre objetivos, afirmaciones y evidencia}

La tabla~\ref{tab:oe-h} relaciona cada objetivo con su proposición analítica o pregunta de evaluación. El OE1 es instrumental.

\begin{table}[H]
\centering
\caption{Correspondencia entre objetivos específicos, afirmaciones, método y evidencia.}
\label{tab:oe-h}
\footnotesize
\setlength{\tabcolsep}{3pt}
\renewcommand{\arraystretch}{1.0}
\begin{tabularx}{\textwidth}{@{}l l >{\raggedright\arraybackslash}X >{\raggedright\arraybackslash}X@{}}
\toprule
Objetivo & Afirmación & Método & Evidencia \\
\midrule
OE1 & Instrumental & Especificación, invariantes y registro de ejecuciones & Nueve verificaciones de invariantes; ID y sumas de verificación (anexo~\ref{cap:anexos}) \\
OE2 & H1, analítica & Paso de capacidad y ticks ideales hasta la banda en forma cerrada; comparación con la grilla simulada & $2\log(1+c)$ por tick aislado; resúmenes simulados de recuperación (sección~\ref{sec:res-h1}) \\
OE3 & H2, exploratoria & Sensibilidad pareada de la regla actual según escala de órdenes y acción de pánico & Contrastes por celda y persistencia de la recuperación (secciones~\ref{sec:res-h2} y~\ref{sec:engineering-sensitivity}) \\
OE4 & H3, comparación de modelos & Descripción predictiva reservada de trayectorias de error frente a referencias simples & Error frente a persistencia y AR(1) escalar; reproducción sensible a la banda (sección~\ref{sec:res-h3}) \\
\bottomrule
\end{tabularx}
\end{table}
